\section{Le dimensionnement à l'épreuve des faits}

Le moteur ne choisit pas seulement \emph{quand} être investi, mais \emph{avec
quelle taille}. Cette taille est censée être pilotée par la volatilité~: plus le
marché s'agite, plus la position se réduit, de façon à viser une volatilité
constante de 55\,\% annualisés. La figure suivante confronte cette intention aux
trente-cinq décisions réellement prises.

\begin{figure}[H]
\centering
\includegraphics[width=\linewidth]{gold_trades_sizing}
\caption{En haut, le levier retenu à l'entrée de chaque trade et son plafond.
En bas, la volatilité réalisée de l'or sur 21~séances, annualisée, et la cible
visée.}
\end{figure}

\begin{insightbox}
\textbf{La cible de volatilité n'est jamais atteinte, et ne peut pas l'être.}
Sur toute la période, la volatilité réalisée de l'or évolue entre 8 et 33\,\%,
c'est-à-dire toujours \emph{sous} la cible de 55\,\%. Le mécanisme demande donc
en permanence plus de levier qu'il n'en obtient~: le levier médian s'établit à
5,2$\times$ et bute onze fois sur trente-cinq contre le plafond de 6,6$\times$.
\end{insightbox}

Cette observation change la façon de lire le paramètre. La «~cible de
volatilité à 55\,\%~» ne fonctionne pas comme un objectif que le moteur
atteindrait~: elle fonctionne comme un \textbf{multiplicateur}, et c'est le
plafond de 6,6$\times$ qui exerce la contrainte effective la plupart du temps.
Régler la cible à la hausse n'aurait presque aucun effet tant que le plafond ne
bouge pas~; c'est bien le couple des deux paramètres qui détermine le risque
porté, et il a été calibré comme tel.

On voit également le mécanisme fonctionner correctement là où il compte. En
avril~2026, la volatilité de l'or dépasse 30\,\% --- la seule période de la
fenêtre où elle s'approche de la cible --- et le levier tombe sous 2$\times$ sur
les quatre trades concernés. Le dimensionnement a donc bien réduit l'exposition
dans le seul épisode d'agitation marquée, ce qui est précisément ce qu'on lui
demande.

\subsection{Ce que le pas de lot impose}

Le simulateur ne peut pas exprimer un levier arbitraire~: il achète des lots
par pas de 0,01, soit une once d'or par cran. Sur les volumes en jeu ici ---
entre 0,09 et 0,26~lot --- ce pas représente entre 3,8 et 11,1\,\% de la
position. L'écart entre le levier voulu et le levier effectivement obtenu est
donc structurellement non nul, et il croît quand le capital diminue.

Ce poste avait été identifié comme un candidat sérieux pour expliquer l'écart de
repli entre les deux moteurs. Les artefacts disponibles ne permettent pas de le
chiffrer --- il faudrait la trace du levier voulu avant arrondi, que le
simulateur ne produit pas. Il reste donc \textbf{non attribué}, et il serait
malhonnête de le présenter autrement.

\begin{warningbox}
Le dimensionnement retenu vise environ 40\,\% de croissance annualisée et a été
calibré sous ce mandat, le risque ayant été déclaré non contraignant. Ce n'est
\textbf{pas} le point qui maximise le rendement~: au-delà du levier optimal, la
traînée de variance l'emporte et le rendement rechute. Augmenter encore la
cible achèterait du repli, pas du rendement.
\end{warningbox}
